%% ****** Start of file apstemplate.tex ****** %
%%
%%
%%   This file is part of the APS files in the REVTeX 4.2 distribution.%%
%%   Copyright (c) 2024 The American Physical Society.
%%
%%   See the REVTeX 4 README file for restrictions and more information.
%%
%
% This is a template for producing manuscripts for use with REVTEX 4.2
% Copy this file to another name and then work on that file.
% That way, you always have this original template file to use.
%
% Group addresses by affiliation; use superscriptaddress for long
% author lists, or if there are many overlapping affiliations.
%  N.B. The groupedaddress option will reorder the author list based
%  on the order in which affiliations appear. Please be sure to check the author 
%  order. You can also use the unsortedaddress(?) option instead to prevent that
%  behavior.
% For Phys. Rev. appearance, change preprint to twocolumn.
% Choose physrev, prl, or rmp for journal
%  N.B. physrev is appropriate for all APS journals except prl and rmp
%  Add 'draft' option to mark overfull boxes with black boxes
%  Add 'showkeys' option to make keywords appear
\documentclass[aps,physrev,preprint,groupedaddress]{revtex4-2}
%\documentclass[aps,physrev,preprint,superscriptaddress]{revtex4-2}
%\documentclass[aps,prl,preprint,superscriptaddress]{revtex4-2}
%\documentclass[aps,prl,reprint,groupedaddress]{revtex4-2}
%\documentclass[aps,rmp,preprint,superscriptaddress]{revtex4-2}
%\documentclass[aps,rmp,reprint,groupedaddress]{revtex4-2}
\usepackage{graphicx}
\usepackage[top=80pt,bottom=100pt,left=68pt,right=66pt]{geometry}
\usepackage{setspace}
\usepackage{amsmath}
\usepackage{amssymb}
\usepackage{hyperref}
%\usepackage{cite}
\usepackage{tablefootnote}
\usepackage{datetime2}
\date{\today} 
\usepackage{changepage}
\usepackage{siunitx}
\usepackage{algorithm}
\usepackage{algpseudocode}
\usepackage[stable]{footmisc}
\usepackage{epigraph}
\usepackage{xcolor}
\newcommand{\ccc}[1]{#1}
\newcommand{\fff}[1]{#1}
\newcommand{\hhh}[1]{#1}
% You should use BibTeX and apsrev.bst for references
% Choosing a journal automatically selects the correct APS
% BibTeX style file (bst file), so only uncomment the line
% below if necessary.
%\bibliographystyle{apsrev4-2}
\newcommand{\kpt}{\mathbf{k}}
\newcommand{\kpp}{\boldsymbol{\kappa}}
\newcommand{\tiKS}{|\mu_{\mathbf{k}n}\rangle}
\newcommand{\tdKS}{|v_{\mathbf{k}n}\rangle}
\newcommand{\tdFL}{|\phi_{\kpt n}\rangle}
\newcommand{\tiFL}{|\varphi_{\kpt n}^\eta\rangle}
\newcommand{\dFKS}{\widetilde{d}_{\kpt ni}(\eta) }
\newcommand{\dFKSj}{\widetilde{d}_{\kpt nj}(\gamma) }
\newcommand{\kplus}{\mathbf{k}_\alpha^\sigma}
\newcommand{\kgrid}[3]{#1$\times$#2$\times$#3}
\newcommand{\xfft}{\mathbf{x}_\text{R}}
\newcommand{\SIfl}[0]{\ref{Ap2}}
\newcommand{\be}{\begin{equation}}
\newcommand{\ee}{\end{equation}}
\newcommand{\ben}{\begin{equation*}}
\newcommand{\een}{\end{equation*}}
\newcommand{\bea}{\begin{eqnarray}}
\newcommand{\eea}{\end{eqnarray}}
\newcommand{\bean}{\begin{eqnarray*}}
\newcommand{\eean}{\end{eqnarray*}}
\newcommand{\lv}{{\bf a}}
\newcommand{\bb}{{\bf b}}
\newcommand{\uu}{{\bf u}}
\newcommand{\qq}{{\bf q}}
\newcommand{\hh}{{\bf h}}
\newcommand{\rr}{{\bf r}}
\newcommand{\pp}{{\bf p}}
\newcommand{\PP}{{\bf P}}
\newcommand{\kk}{{\bf k}}
\newcommand{\HH}{{\bf H}}
\newcommand{\GG}{{\bf G}}
\newcommand{\SiS}{{\bf \Sigma}}
\newcommand{\VV}{{\bf V}}
\newcommand{\UU}{{\bf U}}
\newcommand{\w}{\omega}
\newcommand{\tf}{\textbf}
\newcommand{\bo}{\mathbf}
\newcommand{\br}{{\bf r}}
\newcommand{\nup}{n_{\uparrow}}
\newcommand{\ndown}{n_{\downarrow}}
\newcommand{\Id}[1] {\int \! \! {\rm d}^3 #1}
\renewcommand{\v}[1]{{\bf #1}}
\renewcommand{\[}{\left[}
\renewcommand{\]}{\right]}
\renewcommand{\(}{\left(}
\renewcommand{\)}{\right)}
\def\efield{\boldsymbol{\cal E}}
\def\ket#1{\vert#1\rangle}
\def\bra#1{\langle#1\vert}
\def\susc#1{\chi^{(#1)}}
\def\ket#1{\vert#1\rangle}
\def\bra#1{\langle#1\vert}
\def\ai{\emph{ab-initio}\ }


\begin{document}

% Use the \preprint command to place your local institutional report
% number in the upper righthand corner of the title page in preprint mode.
% Multiple \preprint commands are allowed.
% Use the 'preprintnumbers' class option to override journal defaults
% to display numbers if necessary
%\preprint{}

%Title of paper
\title{Floquet analysis of real-time simulations.}

% repeat the \author .. \affiliation  etc. as needed
% \email, \thanks, \homepage, \altaffiliation all apply to the current
% author. Explanatory text should go in the []'s, actual e-mail
% address or url should go in the {}'s for \email and \homepage.
% Please use the appropriate macro foreach each type of information

% \affiliation command applies to all authors since the last
% \affiliation command. The \affiliation command should follow the
% other information
% \affiliation can be followed by \email, \homepage, \thanks as well.
\author{Claudio Attaccalite}
\author{Ignacio M. Alliati}
%\email[]{Your e-mail address}
%\homepage[]{Your web page}
%\thanks{}
%\altaffiliation{}
\affiliation{}
\author{Myrta Gr\"uning}
%\email[]{Your e-mail address}
%\homepage[]{Your web page}
%\thanks{}
%\altaffiliation{}
\affiliation{}


%Collaboration name if desired (requires use of superscriptaddress
%option in \documentclass). \noaffiliation is required (may also be
%used with the \author command).
%\collaboration can be followed by \email, \homepage, \thanks as well.
%\collaboration{}
%\noaffiliation

\date{\today}

\begin{abstract}
We present a Floquet analysis tool capable of calculating Floquet eigenvectors provided that the time-dependent states, i.e., the solutions of the real-time approach, are known. 
\end{abstract}

% insert suggested keywords - APS authors don't need to do this
%\keywords{}

%\maketitle must follow title, authors, abstract, and keywords
\maketitle

% body of paper here - Use proper section commands
% References should be done using the \cite, \ref, and \label commands
\section{Introduction}
Floquet theory offers a suitable framework to address the time evolution of so-called periodically-driven quantum systems in light of the additional symmetry given by time-periodicity, as initially proposed in the seminal works of Shirley \cite{Shirley1965}, Sambe \cite{Sambe1973} and Salzman \cite{Salzman1973}. This relies on the fact that, provided a time-periodic Hamiltonian, the corresponding time-dependent Schr\"odinger equation is a differential equation with periodic coefficients,
\begin{equation}\label{eq:2-Floquet_1}
    (\hat{\cal H} - i \partial_t)\;\psi(\mathbf{r},t) = 0,
\end{equation}
where $\hat{\cal H}(t+T)=\hat{\cal H}(t)$. In general we would like to condier and Hamiltonian with a time depedence from the external field in the form $\hat{\cal H}(\rho(t+T),t+T)=\hat{\cal H}(\rho(t),t)$, similar to TDDFT or when density matrix is considered to TD-aGW. A differential equation of these characteristics admits the application of Floquet's theorem \cite{Floquet1883}
\footnote{Floquet theory has been covered in numerous textbooks, e.g.,\cite{tannor2007,joachain2011}}, which states that its solutions would be of the form,
\begin{equation}\label{eq:2-Floquet_2}
    \psi(\mathbf{r},t) = e^{-i\xi t}\; \phi(\mathbf{r},t),
\end{equation}
i.e., a pure phase times a periodic function that retains the periodicity of the Hamiltonian, $\phi(\mathbf{r},t+T)=\phi(\mathbf{r},t)$. The latter represents a so-called Floquet state while the argument in the pure phase contains the Floquet quasi-energy, $\xi$. 
% Put \label in argument of \section for cross-referencing
%\section{\label{}}

\section{\label{sec:theo} Theory and methods}
\begin{enumerate}
    \item introducing the real-time
    \item introducing the composite Floquet--Kohn-Sham space (This composite space was originally introduced by Sambe as the Hilbert space for Floquet \textit{steady states} and the functions $e^{-i\fff{\eta\omega_0} t}\mu(\mathbf{r})$ form an (infinite) orthonormal basis for it \cite{Sambe1973}). 
    \item introducing the analysis tool
    \item introduce what we want to analyse and what info we can get
\end{enumerate}


\subsection{\label{sec:realtime} Real-time simulations}
In real-time simulation we evolve valence bands with the equation:
\begin{equation}\label{eq:EOM}
    - i \partial_t  \tdKS =\hat{\cal H} \tdKS
\end{equation}
that are expanded in a Kohn-Sham basis as:
\begin{equation}
	|v_{\mathbf{k}n}(t)\rangle = \sum_{i=1}^{\infty} c_i^n(t) |u_{\mathbf{k}i}\rangle
\end{equation}
The sum is truncated when convergence is achieved in the spectrum.

\subsection{\label{sec:FloqKS} Floquet--Kohn-Sham basis}
Given their periodicity, Floquet states $\phi$ in Eq.~\ref{eq:2-Floquet_2} can be expanded in a Fourier series,
\begin{equation}\label{eq:2-Floquet_3}
    \phi(\mathbf{r},t) = \sum_{\gamma=-\infty}^{+\infty}{\;e^{-i\gamma\omega_0t}\;\widetilde\phi_\gamma(\mathbf{r})},
\end{equation}
where the coefficients, $\widetilde\phi_\gamma(\mathbf{r})$, receive the name of time-independent Floquet states, $\gamma$ runs over the Floquet modes and $\omega_0=\frac{2\pi}{T}$ is the frequency of the perturbation in $\hat{\cal H}$. 

%These exponentials form an orthonormal basis for the Hilbert space of periodic functions of time, ${\cal L}_2[0,T]$, with the inner product,
%\begin{equation}\label{eq:2-Floquet_4}
%    \frac{1}{T}\int_0^T{dt\;e^{+i\eta\omega_0t}\;e^{-i\gamma\omega_0t}} = \delta_{\eta,\gamma}.
%\end{equation}

We further expand the time-independent Floquet states, $\widetilde\phi_\gamma(\mathbf{r})$, in the basis of eigenstates $u(\mathbf{r})$ of the unperturbed Hamiltonian, $\hat H^0$. 
%These functions form an orthonormal set for the Hilbert space we denote ${\cal R}$. Therefore, the eigenproblem in Eq. \ref{eq:2-Floquet_5a} is solved in a composite Hilbert space, ${\cal L}_2[0,T]\otimes{\cal R}$, with the inner product,
%\begin{equation}
%    \langle\hspace{-1.37pt}\langle\cdot|\cdot\rangle\hspace{-1.37pt}\rangle\equiv
%    \frac{1}{T}\int_0^T{dt\; \langle\cdot|\cdot\rangle}, 
%\end{equation}
%where $\langle\cdot|\cdot\rangle$ is the usual inner product in $\mathcal{R}$.


%Another interesting feature of this Floquet formulation relates to the fact that the quasi-energies are defined modulo $\omega_0$. In essence, if $\phi_0$ is an eigenstate of the quasi-energy operator with eigenvalue $\xi_0$, then $\phi_n=e^{-in\omega_0t}\phi_0$ is also an eigenstate with eigenvalue $\xi_n=\xi_0-n\omega_0$. Therefore, $\phi_0$ and $\phi_n$ lead to the same physical state, $\psi$, in Eq. \ref{eq:2-Floquet_2}. This is equivalent to the folding of reciprocal space into the BZ for periodic systems. Indeed, we can define any interval of length $\omega_0$ in the real axis as an $\omega$-BZ, e.g., $[-\omega_0/2,\omega_0/2]$.

\subsection{\label{sec:FloqAn} Floquet analysis tool}
As explained in Section \ref{sec:realtime}, the real-time approach integrates Eq. \ref{eq:EOM} and calculates the time-dependent Bloch states in the KS basis, i.e., $\langle u_{\mathbf{k}i}\tdKS$. For periodically driven systems and under the adiabatic approximation introduced in Eq. \ref{adiabatic_approx}, this magnitude is equal to the projection of Eq. \ref{td-coefficients} over the KS basis,

\begin{equation}
     \label{eq:5-py1}
     \langle u_{\mathbf{k}i} \tdKS  = \sum_{j}{e^{-i\xi_{\kpt n} t} \; d_{\kpt nj}(t)\;\delta_{i,j}} = e^{-i\xi_{\kpt n} t} \; d_{\kpt ni}(t) ,
\end{equation}
where the coefficients $d_{\kpt ni}(t)$ retain the periodicity of the perturbation,
\begin{equation}\label{eq:5-py1b}
    d_{\kpt ni}(t)=d_{\kpt ni}(t+T).
\end{equation}
The idea is to take the output of a real-time calculation, $\langle u_{\mathbf{k}i} \tdKS$, solve Eq.~\ref{eq:5-py1} for $d_{\kpt ni}(t)$ and Fourier transform the latter \textit{via} Eq.~\ref{time-to-eta} to obtain the Floquet eigenvectors, $\dFKS$. However, the quasi-energies $\xi_{\kpt n}$ are not computed explicitly by the real-time approach. In principle, one could obtain the phase from a pair of real-time results separated in time by one period,
\begin{equation}
     \label{eq:5-py2}
     \frac{\langle \mu_{\mathbf{k}i} \tdKS(t_s+T)}{\langle \mu_{\mathbf{k}i} \tdKS(t_s)} = e^{-i\xi_{\kpt n} T}.
\end{equation}
In order to solve the above equation we use a  time step to be commensurate with the period of the field.

\bibliography{books,papers}

\end{document}
%
% ****** End of file apstemplate.tex ******

